\section{Introduction}
Optical music recognition converts score images into machine-readable notation. Unlike ordinary character recognition, the meaning of a mark depends jointly on its class, staff-relative position, attachment relations, voice, and reading order \cite{calvozaragoza2021understanding,rebelo2012optical}. Modern systems increasingly treat OMR as image-to-sequence prediction, using Transformer decoders to emit compact musical encodings \cite{vaswani2017attention,riosvila2024smt,li2023tromr}. This formulation has enabled recognition beyond monophonic incipits and toward piano systems or complete pages \cite{riosvila2023grandstaff,riosvila2026fullpage}. Yet the apparent progress hides a deployment problem: a model trained on one renderer, score collection, vocabulary, or page layout can fail sharply on another.

The shift is visible before any model is run. Synthetic piano pages have clean backgrounds, regular spacing, and constrained engraving, whereas scanned historical scores contain uneven contrast, page texture, warped systems, typography changes, and more varied annotations. The output space shifts as well: Humdrum-based encodings, Linearized MusicXML (LMX), tuplets, voices, and staff markers differ in vocabulary and structural assumptions. General transfer-learning results predict that higher model layers become increasingly task-specific \cite{pan2010transfer,yosinski2014transferable}; for OMR, the decoder and notation vocabulary are precisely the most domain-specific components.

This paper introduces RIGOR-OMR, short for \emph{Relation-aware Interpretable Geometry for Out-of-domain Recognition}. Rather than relying exclusively on a source-specific autoregressive vocabulary, RIGOR-OMR exposes a modular recognition path: staff geometry establishes a normalized coordinate frame; DEIM-D-FINE detects symbols; optional SAM2 masks refine local shape evidence; explicit relations connect noteheads, stems, and beams; and a deterministic backend exports MusicXML or an evaluation-compatible sequence. The intermediate artifacts make it possible to distinguish detection, pitch, duration, relation, and serialization failures.

Our contributions are fourfold:
\begin{itemize}
  \item We formulate practical OMR applicability as zero-shot transfer under matched inputs, frozen target-domain decisions, and a common content projection, rather than extrapolating from in-domain accuracy.
  \item We introduce RIGOR-OMR, an inspectable geometry--detection--relation--reconstruction pipeline whose detector can be replaced without changing the structured backend.
  \item Across Polish and Debussy, local runs against SMT and Zeus show that RIGOR-OMR is significantly stronger than both baselines on Polish and significantly stronger than SMT while statistically comparable to Zeus on Debussy.
  \item Full-factorial and edge-level ablations isolate explicit relation reasoning as the main transferable mechanism, while strict OLiMPiC metrics identify global ordering and hierarchy as the remaining limitation.
\end{itemize}
